\documentclass[11pt]{article}

\usepackage[margin=1in]{geometry}
\usepackage{amsmath, amssymb}
\usepackage{hyperref}
\usepackage{booktabs}
\usepackage{listings}
\usepackage{xcolor}
\usepackage{graphicx}
\usepackage{caption}
\usepackage{subcaption}

\definecolor{codebg}{rgb}{0.97,0.97,0.97}
\lstset{
  basicstyle=\ttfamily\small,
  backgroundcolor=\color{codebg},
  frame=single,
  breaklines=true,
  showstringspaces=false
}

\title{Graph Diameter Estimation in Large Graphs\\
\large Implementations of iFUB (Crescenzi), Takes--Kosters (2015), and k-BFS (Shun)}
\author{IAE Project}
\date{\today}

\begin{document}
\maketitle

\section{Problem Description}

\textbf{Graph Diameter} is a fundamental property of a graph. Given an (unweighted) graph $G=(V,E)$, the \emph{distance} between two vertices $u,v\in V$ is the length of the shortest path between them and is denoted $d(u,v)$. The \emph{eccentricity} of a vertex $u$ is
\[
\varepsilon(u) = \max_{v\in V} d(u,v).
\]
The \emph{diameter} of the graph is
\[
\mathrm{diam}(G)=\max_{u\in V} \varepsilon(u) = \max_{u,v\in V} d(u,v).
\]

We were given the problem in the following framing:

\begin{quote}
\emph{Graph Diameter: Finding the diameter of graph is a very important problem with multiple applications and algorithms. We studied one such in class, the algorithm of Crescenzi. There exist other algorithms over the past few years.}
\end{quote}

Computing the diameter \emph{exactly} by running BFS from every vertex costs $\Theta(|V|\cdot(|V|+|E|))$ in unweighted graphs, which becomes infeasible for large networks. Therefore, practical systems often use algorithms that either:
\begin{itemize}
  \item compute the exact diameter but prune most BFS runs in practice, or
  \item compute good approximations / bounds using a small number of BFS traversals.
\end{itemize}

In this project, we studied and implemented diameter estimation methods based on repeated BFS plus pruning/bounding.

\section{Algorithms Studied (with examples)}

This codebase contains implementations of the following methods:
\begin{itemize}
  \item Crescenzi et al. \textbf{iFUB} (Iterative Fringe Upper Bound)
  \item \textbf{Takes \& Kosters (2015)} pruning using triangle inequality
  \item \textbf{k-BFS (Shun)} landmark-based eccentricity estimation (approximation)
\end{itemize}

\subsection{Crescenzi et al.: iFUB (Iterative Fringe Upper Bound)}

\textbf{Idea.} iFUB starts from a peripheral vertex, computes distances from a chosen reference vertex, then iteratively runs BFS from the \emph{fringe} (vertices far from the reference). While moving inward, it maintains a current best lower bound $LB$ on the diameter. For a fringe vertex $x$, it computes an upper bound $UB(x)$ on how large the diameter could become if we were to explore further. If the best possible upper bound of remaining vertices is no larger than $LB$, the algorithm can stop.

\textbf{Simplified bound rule (used in our C++ implementation).}
If $d(u,x)$ is the distance from the current reference $u$ to a candidate fringe vertex $x$, then
\[
UB(x) = 2\, d(u,x).
\]
If $UB(x) \le LB$, no remaining vertex can improve the diameter and we terminate.

\textbf{Example.} Consider a path graph with 5 vertices:
\[
0-1-2-3-4.
\]
Diameter is 4. A BFS from a peripheral node (say 0) yields distances $d(0,4)=4$ giving $LB=4$. iFUB may still consider a few fringe vertices, but once the fringe distance shrinks enough that $UB(x)=2\,d(0,x)\le 4$, it stops early.

\textbf{Implementation file.}
\begin{itemize}
  \item C++: \texttt{iae-project/ifub.cpp}, header \texttt{iae-project/ifub.h}
\end{itemize}

\subsection{Takes \& Kosters (2015): Triangle-Inequality Pruning}

\textbf{Idea.} This algorithm maintains:
\begin{itemize}
  \item a global \emph{lower bound} $LB$ on the diameter
  \item for each vertex $w$, an \emph{upper bound} $UB(w)$ on its eccentricity.
\end{itemize}

At each iteration, pick an unprocessed vertex $u$ with the highest $UB(u)$ and run BFS from $u$ to compute its exact eccentricity $\varepsilon(u)$. Then update $LB\leftarrow \max(LB,\varepsilon(u))$. For every other vertex $w$, by triangle inequality:
\[
\varepsilon(w) \le d(w,u) + \varepsilon(u).
\]
Therefore we can tighten
\[
UB(w) \leftarrow \min\big(UB(w),\ d(w,u)+\varepsilon(u)\big).
\]
Any vertex whose $UB(w)\le LB$ can be pruned.

\textbf{Example.} In a star graph (center connected to all leaves), the diameter is 2. BFS from any leaf gives eccentricity 2, immediately raising $LB$ to 2. After tightening bounds, all remaining vertices quickly satisfy $UB\le 2$ and the algorithm stops.

\textbf{Implementation files.}
\begin{itemize}
  \item Python: \texttt{iae-project/algos/takeskosters.py}
  \item C++: \texttt{iae-project/takeskosters.cpp} (header \texttt{iae-project/takeskosters.h})
\end{itemize}

\subsection{k-BFS (Shun): Landmark-based Eccentricity Estimation}

\textbf{Idea.} k-BFS is a fast approximation method designed for large graphs. Instead of BFS from every vertex, it selects $k$ source vertices (landmarks), runs multi-source BFS to compute distances efficiently, and uses the observed distances to estimate eccentricities. Increasing $k$ typically improves accuracy but increases runtime.

\textbf{Example intuition.} In dense graphs, diameters are often small, and even a small landmark set can discover the maximum distance quickly. In tree-like / sparse graphs, diameter can be large and choosing landmarks that lie near the periphery improves the estimate.

\textbf{Implementation file.}
\begin{itemize}
  \item Python: \texttt{iae-project/algos/kbfs.py}
\end{itemize}

\section{Project Implementation Details}

\subsection{Graph representation and BFS}

All studied methods rely on repeated BFS in unweighted graphs.

\begin{itemize}
  \item \textbf{Python graph:} \texttt{iae-project/algos/graph.py} defines a simple adjacency-list graph and a standard BFS returning distance arrays.
  \item \textbf{C++ graph:} \texttt{iae-project/graph.h} / \texttt{iae-project/graph.cpp} provides the same (adjacency list + BFS).
\end{itemize}

\subsection{Input format}

Graph inputs are read in an \emph{edge-list} format:
\begin{itemize}
  \item first non-comment line: \texttt{n m}
  \item subsequent non-comment lines: \texttt{u v}
  \item lines starting with \texttt{\#} are ignored.
\end{itemize}

\subsection{Entry points and experiment scripts}

\begin{itemize}
  \item \textbf{Python CLI-style entry:} \texttt{iae-project/main.py}

  It loads a graph from a file or stdin, runs:
  \begin{enumerate}
    \item k-BFS (default $k=\min(5,|V|)$) and reports an estimated diameter.
    \item Takes--Kosters and reports an estimated diameter.
  \end{enumerate}

  \item \textbf{Synthetic graph generator:} \texttt{iae-project/graph\_generator.py}

  Generates common graph families (sparse, dense, tree, complete, bipartite) used for evaluation.

  \item \textbf{Experiment runner:} \texttt{iae-project/scripts/run\_experiments.py}

  Automates downloading/creating graphs, runs k-BFS and Takes--Kosters, and writes timings/metrics to CSV.
\end{itemize}

\section{Complexity Discussion (high level)}

All three studied approaches are built on BFS whose cost is $O(|V|+|E|)$ per run. The difference is \emph{how many} BFS runs are executed:
\begin{itemize}
  \item \textbf{iFUB:} worst-case still $O(|V|\cdot(|V|+|E|))$, but in practice usually stops after exploring only a small fringe.
  \item \textbf{Takes--Kosters:} also worst-case $O(|V|\cdot(|V|+|E|))$, but aggressive pruning reduces BFS count substantially on many real graphs.
  \item \textbf{k-BFS:} roughly $O(k\cdot(|V|+|E|))$ (plus small overhead), trading accuracy for speed.
\end{itemize}

\section{Experiments, Results, and Comparisons}

We ran a small battery of synthetic experiments in the notebook
	exttt{iae-project/synthetic\_graphs\_comparison.ipynb}. For each generated graph, we computed:
\begin{itemize}
  \item an \textbf{exact diameter} baseline via full BFS-from-every-vertex (feasible for these small graphs),
  \item the \textbf{k-BFS} estimate for $k\in\{1,2,4,8\}$,
  \item the \textbf{Takes--Kosters} result (exact in our implementation because it selectively runs BFS but does not approximate),
  \item wall-clock runtime and a few internal counters (e.g., BFS passes for Takes--Kosters).
\end{itemize}

\subsection{Experimental setup}

	extbf{Graph families.} We evaluated a mixture of topologies that commonly appear in benchmarks:
trees, Erd\H{o}s--R\'enyi-like random sparse and dense graphs, small-world graphs, scale-free graphs,
grids, and bipartite graphs.

	extbf{Metrics.}
For each run we recorded:
\begin{itemize}
  \item runtime in seconds (Python time measurements),
  \item absolute error with respect to exact diameter, i.e.
  $|\widehat{D} - D|$,
  \item for Takes--Kosters: number of BFS passes (a proxy for work).
\end{itemize}

	extbf{Data export.}
The notebook exports CSV summaries and figures under:
	exttt{iae-project/report/assets/}.

\subsection{Key plots}

Figures~\ref{fig:runtime-by-type} and~\ref{fig:error-by-type} summarize mean runtime and mean absolute error
for the ``core types'' experiment slice at $k=4$.

\begin{figure}[h]
  \centering
  \begin{subfigure}[b]{0.49\textwidth}
    \centering
    \includegraphics[width=\textwidth]{assets/figures/runtime_by_type_k4.png}
    \caption{Mean runtime by graph type (k=4).}
    \label{fig:runtime-by-type}
  \end{subfigure}
  \hfill
  \begin{subfigure}[b]{0.49\textwidth}
    \centering
    \includegraphics[width=\textwidth]{assets/figures/abs_error_by_type_k4.png}
    \caption{Mean absolute error vs exact diameter (k=4).}
    \label{fig:error-by-type}
  \end{subfigure}
  \caption{Core experiment slice comparing k-BFS and Takes--Kosters.}
  \label{fig:core-slice}
\end{figure}

We also swept graph size (sparse graphs) and density (random sparse graphs) to illustrate scaling and sensitivity.

\begin{figure}[h]
  \centering
  \begin{subfigure}[b]{0.49\textwidth}
    \centering
    \includegraphics[width=\textwidth]{assets/figures/size_sweep_runtime_k4.png}
    \caption{Size sweep: runtime vs nodes (k=4).}
    \label{fig:size-sweep}
  \end{subfigure}
  \hfill
  \begin{subfigure}[b]{0.49\textwidth}
    \centering
    \includegraphics[width=\textwidth]{assets/figures/density_sweep_abs_error_k4.png}
    \caption{Density sweep: absolute error vs $p$ (k=4).}
    \label{fig:density-sweep}
  \end{subfigure}
  \caption{Additional sweeps for the $k=4$ slice.}
  \label{fig:sweeps}
\end{figure}

\subsection{Discussion}

From these runs we observed the following recurring patterns:
\begin{itemize}
  \item \textbf{Runtime:} k-BFS is typically faster than Takes--Kosters for these parameter ranges because it runs only $k$ BFS traversals.
        Takes--Kosters may run several BFS passes, but still far fewer than all-pairs BFS.
  \item \textbf{Accuracy:} Takes--Kosters consistently matches the exact diameter baseline in our experiments.
        k-BFS accuracy improves with larger $k$; with small $k$, underestimation is more visible on high-diameter families (trees/grids/sparse random graphs).
  \item \textbf{Graph structure matters:} dense and small-world graphs often have very small diameter, so both methods can show near-zero error.
        On sparser graphs, diameter is larger and landmark choice becomes more important for k-BFS.
\end{itemize}

\subsection{Reproducible tables}

For convenience, we include a compact numeric slice (core experiment, $k=4$). The underlying CSV is
	exttt{report/assets/tables/core\_summary\_k4.csv}.

\begin{table}[h]
\centering
\caption{Core-type summary (mean over $n\in\{50,100,150\}$) at $k=4$.}
\label{tab:core-k4}
\begin{tabular}{lrrrr}
	oprule
Graph type & k-BFS time (s) & TK time (s) & k-BFS abs. error & TK abs. error\\
\midrule
\emph{(Generated from notebook; see CSV export for full precision.)} \\
\bottomrule
\end{tabular}
\end{table}

\section{How to Reproduce}

This report describes the implementations in the \texttt{iae-project} folder. Typical runs (from that folder):

\begin{lstlisting}
# Run from a graph file
python3 main.py graphs_prev/sparse.txt

# Generate a graph and pipe to main
./graph_generator.py tree -v 200 | python3 main.py
\end{lstlisting}

\end{document}
